%! AP Statistics — Unit 2, Lesson 2.8 — Homework
\documentclass[10pt]{article}
\usepackage{apstats-article}
\usepackage{apstats-boxes}

\begin{document}

\pageheader{Unit 2, Lesson 2.8}{Homework}
\namedateperiod

\vspace{0.14in}

% ── PROBLEM 1 ─────────────────────────────────────────────────────────────────
\begin{scenariobox}[Context A: Study Hours and GPA]{sky}
\small A researcher records weekly study hours ($x$) and semester GPA ($y$) for 50 college
students. Summary statistics: $\bar{x} = 18$ hr, $\bar{y} = 3.1$, $s_x = 6$ hr, $s_y = 0.5$,
$r = 0.70$. The data ranged from 5 to 35 hours per week.
\end{scenariobox}

\vspace{0.1in}

\begin{enumerate}[leftmargin=1.6em, itemsep=0.18in]

  \item Compute $b$ and $a$. Show work.\\[8pt]
        \writeline\\[3pt]\writeline\\[3pt]\writeline

  \item Write the regression equation using variable names.\\[8pt]
        \writeline

  \item Interpret the \textbf{slope} in context.\\[8pt]
        \writeline\\[3pt]\writeline

  \item Interpret the \textbf{$y$-intercept} in context. Is it meaningful? Explain.\\[8pt]
        \writeline\\[3pt]\writeline\\[3pt]\writeline

  \item Compute $r^2$. Write the full AP-template interpretation.\\[8pt]
        \writeline\\[3pt]\writeline\\[3pt]\writeline

  \item Predict the GPA for a student who studies 25 hours per week. Classify as interpolation
        or extrapolation.\\[8pt]
        \writeline\\[3pt]\writeline

\end{enumerate}

\vspace{0.18in}

% ── PROBLEM 2 ─────────────────────────────────────────────────────────────────
\begin{scenariobox}[Context B: Altitude and Temperature]{goldbg}
\small A meteorologist measures altitude in thousands of feet ($x$) and air temperature in
\textdegree F ($y$) at 40 locations. Summary statistics: $\bar{x} = 5$ (thousand ft),
$\bar{y} = 42$\textdegree F, $s_x = 2$ (thousand ft), $s_y = 14$\textdegree F, $r = -0.90$.
\end{scenariobox}

\vspace{0.1in}

\begin{enumerate}[leftmargin=1.6em, itemsep=0.18in, start=7]

  \item Compute $b$, $a$, and $r^2$. Write the regression equation.\\[8pt]
        \writeline\\[3pt]\writeline\\[3pt]\writeline

  \item Interpret $r^2$ in context.\\[8pt]
        \writeline\\[3pt]\writeline

  \item A data point at altitude 18,000 ft ($x = 18$) with a temperature of $-65$\textdegree F
        is discovered. Is this likely to be an influential point? Why?\\[8pt]
        \writeline\\[3pt]\writeline\\[3pt]\writeline

\end{enumerate}

\vspace{0.18in}

% ── EXTENSION ────────────────────────────────────────────────────────────────
\begin{extensionbox}
\small\textbf{Extension --- optional}

\vspace{4pt}
For Context B, the point $(18, -65)$ is added. The new regression equation becomes
$\widehat{\text{temp}} = 82 - 8.5(\text{altitude})$ and $r^2 = 0.94$.

\begin{enumerate}[leftmargin=1.6em, itemsep=6pt]
  \item How did the slope change? The intercept? The $r^2$? Describe each change.
  \item Does this confirm or disconfirm that the point is influential? Explain.
  \item If you were the meteorologist, would you remove this point? What would you investigate first?
\end{enumerate}
\end{extensionbox}

\vspace{0.16in}

% ── PREVIEW ──────────────────────────────────────────────────────────────────
\begin{tcolorbox}[colback=greenbg, colframe=greenacc, boxrule=0.8pt, arc=2mm,
    left=4mm, right=4mm, top=2.5mm, bottom=2.5mm]
\small\textbf{Preview for Lesson 2.9:} The LSRL is the best \emph{linear} model --- but what
if the relationship is curved? When a residual plot shows a pattern (not random scatter), a
\textbf{transformation} (log, square root, power) can straighten the data before applying
regression. We'll explore when and how to transform variables next class.
\end{tcolorbox}

\end{document}
